\subsection{Relation annotations and source captions}
We use the official e-ViL release of e-SNLI-VE at revision \texttt{94f17f5c1895}. Its training, development, and test sources contain 401,717, 14,339, and 14,740 annotated hypotheses over 29,783, 1,000, and 1,000 images, respectively. The official pools have no shared image IDs or annotation-pair IDs. With Python's \texttt{random.Random(42)}, we sample from sorted image IDs: 1,200 training images, 200 development images split into 100 calibration and 100 validation images, and 400 test images. All annotations of each selected image remain in source order. Table~\ref{tab:split-counts} gives the resulting counts.

% Generated from audited fresh results. Do not edit numeric content.
\begin{table}[t]
\centering
\setlength{\tabcolsep}{2.3pt}
\begin{tabular}{lrrrr}
\toprule
Split & Images & Support & Contradict & Neutral\\
\midrule
Train & 1,200 & 5,319 & 5,846 & 5,027\\
Calibration & 100 & 522 & 554 & 355\\
Validation & 100 & 502 & 534 & 339\\
Test & 400 & 2,088 & 2,278 & 1,495\\
\bottomrule
\end{tabular}
\caption{Retained hypothesis annotations. Each image also has five source captions. All validation and test images, and 1,192 of 1,200 training images, have both supported and contradicted hypotheses.}
\label{tab:split-counts}
\end{table}


The source captions come from the official Flickr30k Entities annotation archive at revision \texttt{68b3d6f12d1d}. We use its five sentence-file lines per image in stored order. Entity markup is removed while retaining source words, case, and punctuation, with single spaces between stored tokens. This source covers all 31,783 images. An incomplete alternative caption mirror covers only 31,014 and was not used. The image bytes come from the \texttt{nlphuji/flickr30k} mirror at revision \texttt{2b239befc81b}. Each selected image passes ZIP size/CRC checks, JPEG verification, and full decoding. No selected image is silently skipped.

An earlier deprecated relation release was rejected during an input-only audit because of development/test overlap. This decision preceded feature-based evaluation. The chosen official release, exact IDs, source URLs, full revisions, file hashes, and selection order are recorded in the per-dataset provenance files. The released labels include inherited text-inference judgments, so their names are operational annotation categories rather than a claim of new visual adjudication.

\subsection{Transfer benchmarks}
SugarCrepe uses all 7,511 pairs over 1,560 images from revision \texttt{0047054b2439}. Its category counts are 692 added attributes, 2,062 added objects, 788 replaced attributes, 1,652 replaced objects, 1,406 replaced relations, 666 swapped attributes, and 245 swapped objects. SugarCrepe++ uses all 4,757 triplets over 1,542 images from revision \texttt{c083bbf039a5}; it covers the five replacement/swap categories with the same category counts.

Each benchmark preserves its own released text strings. SugarCrepe++ changes 1,203 first-positive strings and 878 foil strings relative to the matched SugarCrepe records. Its first-positive accuracy therefore cannot substitute for SugarCrepe's replacement/swap score. Two released SugarCrepe++ triplets have identical first-positive and foil strings, making a strict win impossible; these remain in the official denominator. One triplet repeats the same string in both positive positions. No item is edited or removed in response to model scores.

COCO uses all 5,000 Stanford Karpathy test images. The source contains 25,010 captions; taking the first five in stored order for every image yields 25,000 candidates and omits ten sixth captions. Caption strings are not normalized. The COCO JPEG cache contains 6,377 distinct benchmark images. It uses the official \texttt{val2014} source for Karpathy members and \texttt{val2017} for remaining benchmark images. All 183 shared SugarCrepe/Karpathy images were checked across these source versions and are byte-identical.

\subsection{Overlap and duplicate audit}
All 1,542 SugarCrepe++ images also occur in SugarCrepe. COCO retrieval shares 183 images with SugarCrepe and 182 with SugarCrepe++. These are complementary measurements on overlapping pools. The selected Flickr30k images share neither IDs nor image-byte hashes with the COCO-derived benchmarks. The four e-SNLI-VE splits also have no shared image bytes.

The input audit finds 48 repeated exact image/text groups in e-SNLI-VE, each appearing twice. Of these, 46 are repeated hypothesis groups. Their relation labels do not conflict. All rows are retained, so repeated annotations retain their multiplicity in the loss and metrics. The two positive/foil equality cases account for the within-triplet relation conflicts in SugarCrepe++. The audit records source defects before scoring rather than changing benchmark contents.

\subsection{Frozen feature extraction}
The encoder is \path{laion/CLIP-ViT-B-32-laion2B-s34B-b79K}, revision \texttt{1a25a446712b}, with weight SHA256 beginning \texttt{ac4f8c4b88af}. We use OpenCLIP 2.32.0 and PyTorch 2.7.1 on CPU. Images are resized with bicubic interpolation, center-cropped to $224\times224$, converted to RGB, and normalized with the checkpoint's channel means and standard deviations. The 512-dimensional image and text outputs are stored in float32 after unit L2 normalization. The fixed scale is $100.00000762939453$.

Text encoding trims only causal padding positions after the last end-of-text token in a batch. A parity check compares normalized outputs against the stock encoder with tolerance $2\cdot10^{-6}$. Recorded per-dataset parity checks, preprocessing parameters, source hashes, input-manifest hashes, feature-cache hashes, and full model pins are stored in each feature metadata file. Reused features require exact image-byte or text-string identity under the same encoder and preprocessing.
